% 来源：OEBPS/Text/01_55.xhtml（OPF spine 顺序；仅第一册）
\cleardoublepage
\BookEntry[西塞罗《论老年论友谊论责任》]{西塞罗《论老年论友谊论责任》}

\BookCover{images/书单38论老年论友谊论责任.jpg}{[古罗马] 西塞罗著| 徐奕春译| 商务印书馆2003年版}

\begin{BookSummaryBox}
\textbf{古罗马哲学家、罗马文学黄金时代的天才作家西塞罗畅谈老年和友谊的传世名篇}
\par
\end{BookSummaryBox}

\subsubsection*{这本书谈什么？}

西塞罗在《论老年论友谊论责任》中，从积极参与社会政治生活的市民角度出发，谈论如何看待老年的问题。与以往人们认为老年生活十分悲惨的看法不同，西塞罗大力赞美老年生活。
\par

诚然，退休、体能衰弱、无法享受生理上的快乐，及死亡迫近会给老年人生活带来不一样的体验，但事实上，老人无所谓退休，可以活到老学到老；体能衰弱并不表示他们对国家毫无贡献。其次，西塞罗强调老年人要顺应自然规律，适度地运动，保持活力。再者，生理欲求的下降，反而使得老年人更能专注于智性与德性的成长。而最后，死亡是自然之事，人终将一死，但要能平静以对，必须对死亡有毕生的思考。
\par

\subsubsection*{芒格为什么谈起这本书？}

2006年，芒格的朋友安格斯·麦克拜恩夫妇送给他一本书富兰克林1744年译文重印版的《西塞罗论美好人生》。他很喜欢这部充满对老年的赞颂的作品，并在其中找到了赞同自己生活方式的内容。在西塞罗看来，抱怨年纪是无益的，尤其是年轻人不要未老先死，如果你的生活方式是正确的，那么到了晚年会更幸福。芒格经常引用西塞罗的下面这段话：
\par

\begin{quotation}
晚年的最佳保护铠甲是一段在它之前被悉心度过的生活，一段被用于追求有益的知识、光荣的功绩和高尚的举止的生活；过着这种生活的人从青年时代就致力于提升他自己，而且将会在晚年收获它们产生的最幸福的果实；这不仅是因为有益的知识、光荣的功绩和高尚的举止将会陪伴他终生，甚至直到他生命的最后一刻，也会因为见证了正直的人生的良心和过往美好功绩的回忆将会给灵魂带来无上的安慰。
\par
\end{quotation}

\subsubsection*{《论老年论友谊论责任》作者西塞罗如是说}

“历史是时代的见证，真理的火炬，记忆的生命，生活的老师和古人的使者。”
\par

“如果你对自己出生之前的事情毫不了解，那么你永远都是一个无知的孩童。因为人生如果不是对于历史密不可分的往事的记忆，又能是什么呢？”
\par
